\section{Truncating the high spectrum of a Hermitian matrix}
\label{sec:energy_truncation}

For a Hermitian matrix $A$ and an energy $\tau$, let $P_{\ge\tau}$ be the
spectral projection of $A$ onto eigenvalues at least $\tau$.

\begin{definition}[Energy truncation]
    \label{def:energy_truncation}
    \lean{Matrix.energyTruncation}
    \leanok
    \uses{def:closed_threshold_projection}
    The truncation of $A$ at $\tau$ is $\tilde A=\min(A,\tau)$ in the functional
    calculus: every eigenvalue above $\tau$ is replaced by $\tau$. This is the
    truncated Hamiltonian of Arad, Kuwahara and
    Landau~\cite{Arad2014Connecting}, Definition~def:Ht, and the hard truncation
    of~\cite[lines 873--886]{Arad2016RigorousRG}, Definition~def:hard-trunc,
    with $A=H_J$ and $\tau=\varepsilon_J+t$.
\end{definition}

\begin{theorem}[Properties of the energy truncation]
    \label{thm:energy_truncation}
    \lean{Matrix.IsHermitian.energyTruncation_eq, Matrix.isHermitian_energyTruncation,
        Matrix.IsHermitian.commute_energyTruncation, Matrix.IsHermitian.energyTruncation_le,
        Matrix.IsHermitian.energyTruncation_le_smul_one,
        Matrix.IsHermitian.smul_one_le_energyTruncation}
    \leanok
    \uses{def:energy_truncation}
    The truncation equals $A(I-P_{\ge\tau})+\tau P_{\ge\tau}$, is Hermitian and
    commutes with $A$, and satisfies $\tilde A\le A$ and $\tilde A\le\tau I$. If
    every eigenvalue of $A$ is at least $\varepsilon$ and $\varepsilon\le\tau$,
    then $\varepsilon I\le\tilde A$.
\end{theorem}

\begin{proof}
    \leanok
    All four operators are functions of $A$. The first identity is the
    pointwise identity $\min(t,\tau)=t\,\mathbf 1_{t<\tau}+\tau\,\mathbf
    1_{t\ge\tau}$. The order bounds follow from $\min(t,\tau)\le t$,
    $\min(t,\tau)\le\tau$ and $\varepsilon\le\min(t,\tau)$ on the spectrum, by
    monotonicity of the functional calculus.
\end{proof}
